% Chapter 2 -- Background and Related Work
% Standalone, compilable draft (FPT-template article class + booktabs), sibling to
% sec_3_detection_pipeline.tex / sec_8_discussion.tex. Positions this thesis's
% classical-clustering + binary-classifier + PCN-completion pipeline against the
% deep-learning LiDAR-segmentation and point-completion literature, and carries the
% B3 mandatory literature-positioning table (Table 2.1) that defends examiner Q2
% ("why F1 0.808 / R 0.73 vs SOTA?") and Q10 ("what is yours vs. reimplementation?").
%
% MERGE/NUMBERING NOTE: this file opens with \section{Background and Related Work} for
% standalone compilation and numbering, but at final assembly the assembler promotes it
% to a chapter (Chapter 2). Forward references to sibling chapters use typed numerals
% ("Chapter~\ref{ch:pipeline}", "Chapter~\ref{ch:recall-bottleneck}", "Chapter~\ref{ch:completion-eval}") because the chapter files do not share a label
% namespace when built separately.
%
% CITATIONS: unlike the results/methods siblings (which cite the internal Finding record),
% this is a literature chapter and uses classic BibTeX. \bibliography{../../references}
% resolves to docs/references.bib (same relative-path pattern as
% docs/report/progress_report_2026_05_12.tex, confirmed to compile in this repo).
%
% Scope decision (settled with advisor/author): RELATED-WORK-FOCUSED. Method mechanics
% (PointNet architecture, HDBSCAN internals, PCN encoder/decoder, Chamfer distance) are
% explained in Chapter~\ref{ch:pipeline} (pipeline) and Chapter~\ref{ch:completion-method} (completion) and are NOT re-derived
% here; this chapter carries only the light conceptual grounding needed to read the
% related work and the positioning table.
%
% Honesty-phrasing rules applied (personal_agenda P1):
%   - the positioning table (Table 2.1) presents NO protocol-mismatched number as
%     head-to-head: each cell is labelled with its own metric, requirements columns lead,
%     and the caption states the metric/protocol/scope mismatch explicitly
%   - every quantitative claim about a published method traces to a cited primary source;
%     all Table 2.1 numbers are SemanticKITTI test-set figures verified against the paper
%   - this work's own row uses the frozen headline numbers (F1 0.808 / R 0.730 / P 0.905,
%     seq 08; Chapter~\ref{ch:detection-eval}, Finding~#34) with the point-IoU>=0.3 protocol named
%   - "recall" is stated against its car-only, >=10-surviving-points denominator where used
%
\chapter{Background and Related Work}
\label{ch:background}

This chapter places the thesis in its research context. It reviews the two bodies of
work the pipeline draws on --- LiDAR point-cloud segmentation and object-level
point-cloud completion --- and then positions the specific design choices of this
work against them. The organising contrast runs throughout: the dominant approach to
LiDAR segmentation is a single end-to-end network trained on dense per-point
annotation, whereas this thesis builds an interpretable pipeline from training-free
geometric clustering plus a lightweight learned classifier, and treats completion as a
separate, downstream problem. Section~\ref{sec:bg-positioning} makes that contrast
quantitative in a positioning table. The mechanics of each component are deferred to
the method chapters (Chapter~\ref{ch:pipeline} for detection, Chapter~\ref{ch:completion-method} for completion); here the aim
is to say what exists, and where this work sits within it.

\section{LiDAR point clouds and the SemanticKITTI benchmark}
\label{sec:bg-dataset}

A rotating LiDAR sensor returns, for each scan, an unordered set of 3D points sampled
from the surfaces around the vehicle. The points are sparse, irregularly spaced, and
strongly range-dependent: nearby objects are dense, distant ones may be a handful of
returns, and every object is seen from a single viewpoint and is therefore partial.
These properties --- irregularity, sparsity, single-view partiality --- shape every
method discussed below and every difficulty this thesis reports.

The KITTI Vision Benchmark~\cite{geiger2012we} established the standard automotive
LiDAR setting (a Velodyne HDL-64E on a car, with synchronised camera and pose). Its
segmentation-oriented successor, SemanticKITTI~\cite{Behley_2019_ICCV}, provides
dense per-point semantic labels for the 22 sequences of the KITTI odometry split
(sequences 00--10 are released with labels; the labels for sequences 11--21 are withheld
and scored only through the online test server), and a later extension adds temporally
consistent instance identities for a LiDAR-based panoptic benchmark~\cite{behley2021panoptic}. Three tasks are defined on
it. \emph{Semantic segmentation} assigns every point a class and is scored by the mean
intersection-over-union across classes (mIoU). \emph{Instance segmentation} separates
individual objects within the ``thing'' classes (car, pedestrian, cyclist, and so on).
\emph{Panoptic segmentation} unifies the two and is scored by panoptic quality (PQ),
which factorises into a segmentation-quality term (SQ, the mean IoU of matched
segments) and a recognition-quality term (RQ, an F$_1$-like matching score);
PQ$^{\text{Th}}$ restricts PQ to the ``thing'' classes, and PQ$^\dagger$ is a variant
that scores ``stuff'' by IoU. These definitions matter for
Section~\ref{sec:bg-positioning}, because they are the metrics the literature reports
and none of them is the point-level per-class F$_1$ this thesis uses.

The decisive property of SemanticKITTI for this work is its \emph{annotation cost}.
The labels are dense and per-point across billions of points, and the panoptic
extension adds a manual instance-association pass over sequences. This is the resource
the mainstream methods below consume and the one this thesis is built to avoid: its
segmentation stage requires no labelled point clouds at all, and its classifier is
trained on clusters mined and weakly labelled automatically from the pipeline's own
output (Chapter~\ref{ch:pipeline}).

\section{Deep segmentation of LiDAR point clouds}
\label{sec:bg-deep}

The current benchmark leaders are deep networks trained end-to-end on the dense labels.
They differ mainly in how they impose structure on the irregular point set.
\emph{Projection-based} methods rasterise the scan into a 2D range image and apply
image convolutions: SqueezeSeg~\cite{10.1109/ICRA.2018.8462926} and
RangeNet++~\cite{milioto2019rangenet++} are representative, and are fast because they
reuse mature 2D machinery, at the cost of discretisation artefacts at depth
discontinuities. \emph{Point-based} methods operate on the raw points through
permutation-invariant operators: PointNet~\cite{Qi_2017_CVPR} introduced the
shared-MLP-plus-symmetric-pooling design, PointNet++~\cite{NIPS2017_d8bf84be} added
hierarchical local grouping, and RandLA-Net~\cite{Hu_2020_CVPR} made the point-based
route scale to full driving scans through efficient random sampling. \emph{Voxel-based}
methods discretise 3D space and run sparse convolutions; Cylinder3D~\cite{zhu2021cylindrical}
uses a cylindrical partition matched to the sensor's polar geometry and reports
\num{67.8} mIoU on the SemanticKITTI test set, among the higher single-scan results.

Panoptic methods extend semantic backbones with instance reasoning.
Panoptic-PolarNet~\cite{zhou2021panopticpolarnet} predicts instances proposal-free in a
polar bird's-eye-view; DS-Net~\cite{hong2021dsnet} clusters in a learned embedding via a
dynamic shifting module; Panoster~\cite{gasperini2021panoster} learns instance ids
end-to-end; and EfficientLPS~\cite{sirohi2021efficientlps} targets an efficient
architecture for the joint task. The dataset authors' own baselines pair a semantic
network with a 3D detector to supply instances~\cite{behley2021panoptic}.

Two properties are common to this whole family and frame the thesis's alternative.
First, all of them require the dense per-point (and, for panoptic, per-instance)
supervision described above; their accuracy is inseparable from that annotation budget.
Second, they are end-to-end networks whose decisions are not individually inspectable
--- a misclassified region has no stage-level explanation. PointNet and PointNet++ are
relevant to this thesis for a narrower reason as well: they are the backbone family from
which its car/not-car classifier is built (Chapter~\ref{ch:pipeline}). The classifier, however, is a
small binary network applied to already-segmented clusters, not a full-scene
segmentation model, and it is trained without any hand-drawn labels.

\section{Classical geometric segmentation}
\label{sec:bg-classical}

Before end-to-end learning dominated the benchmarks, LiDAR objects were extracted by
geometric reasoning, and that route is the one this thesis takes for the segmentation
stage. The standard recipe removes the ground plane --- commonly by RANSAC plane
fitting --- and clusters the remaining points into objects. The clustering step is
where the design choice lives. Fixed-radius methods such as Euclidean clustering and
DBSCAN~\cite{ester1996dbscan} group points within a single distance threshold
$\varepsilon$; this is fast and fully interpretable, but a single $\varepsilon$ cannot
serve both the dense near field and the sparse far field of a LiDAR scan.
HDBSCAN~\cite{campello2013hdbscan}, in the widely used implementation
of~\cite{mcinnes2017hdbscan}, replaces the fixed radius with a hierarchy over density
levels and extracts clusters that are stable across scales, which is better matched to
range-varying point density.

This clustering-plus-classification recipe remains an active design point wherever
interpretability and real-time efficiency are valued over benchmark accuracy. Recent
work for autonomous racing, for instance, pairs bird's-eye-view clustering with a
lightweight machine-learning vehicle classifier to reach long-range detection without a
deep dense-feature network~\cite{lim2025longrange}. The pipeline in this thesis is of
the same family --- ground removal, clustering, then a learned per-object classifier ---
applied to the SemanticKITTI setting.

The appeal of this route is that it needs no training data for segmentation and every
stage is inspectable. Its weaknesses are equally concrete and are documented directly in
this thesis: density-based clustering has no notion of objectness, so it fragments a
single car into several clusters when internal density gaps appear, and this
fragmentation --- not classification error --- is the dominant recall limiter
(Chapter~\ref{ch:recall-bottleneck}). The positioning below should be read with that trade in view: the
interpretability and near-zero annotation of the classical route are bought at a
measurable recall cost relative to the supervised networks.

\section{Object-level point-cloud completion}
\label{sec:bg-completion}

The second phase of the pipeline reconstructs a denser shape from the sparse, single-view
partial cloud of a detected car. Point-cloud completion learns this partial-to-complete
mapping. PCN~\cite{yuan2019pcnpointcompletionnetwork} established the coarse-to-fine
design --- a PointNet encoder producing a global shape code, decoded first to a coarse
set and then folded to a dense output --- and remains a strong, compact baseline. Later
work improves detail through attention and progressive refinement:
PoinTr~\cite{yu2021pointr} casts completion as set-to-set translation with a
geometry-aware transformer, SnowflakeNet~\cite{xiang2021snowflakenet} generates points
by hierarchical ``snowflake'' deconvolution, and SeedFormer~\cite{zhou2022seedformer}
introduces patch-seed representations. These models are trained on synthetic CAD shapes,
most commonly ShapeNet~\cite{chang2015shapenetinformationrich3dmodel}, with partial
inputs produced by virtual rendering.

Two points from this literature bear directly on later chapters. First, models trained
on clean synthetic partials face a synthetic-to-real domain gap when applied to noisy,
sparse LiDAR clusters; this thesis adopts PCN and confronts that gap head-on
(Chapter~\ref{ch:completion-method}), and its comparison of PCN against a transformer alternative found no
decisive real-data advantage for the heavier model (Chapter~\ref{ch:completion-eval}). Second, and more
consequentially, the standard completion metric --- Chamfer distance to a reference
shape --- presupposes a complete ground-truth surface, which real accumulated LiDAR does
not provide: the accumulation is itself one-sided. Chapter~\ref{ch:completion-eval} shows this makes pseudo
ground-truth Chamfer scoring invalid on real cars and develops an alternative metric in
its place. This should be distinguished from \emph{semantic scene completion}, which
fills a whole voxelised scene from a single scan; the present work completes individual
segmented objects, a different and narrower task.

\section{Positioning of this work}
\label{sec:bg-positioning}

The pipeline studied here is a car-only instance detector built from training-free
clustering and a binary classifier, followed by object completion. It occupies a
different point in the design space from the benchmark leaders of
Section~\ref{sec:bg-deep}, and the honest way to state that is to show the published
numbers alongside its own while being explicit that they are not the same measurement.
Table~\ref{tab:positioning} does this. The columns that carry the argument are the
requirement columns --- approach and supervision --- not the last column: the reported
figures are shown to be candid about the accuracy gap, not to claim a win on a shared
axis, and no two cells in that column are the same metric unless their headers say so.

\begin{table}[H]
  \centering
  \caption{Positioning against representative SemanticKITTI methods. \textbf{The
    figures in the last column are not comparable across rows and must not be read as a
    ranking.} They are different metrics (semantic mIoU and per-class IoU; panoptic PQ,
    PQ$^{\text{Th}}$, PQ$^\dagger$; and this work's point-level car F$_1$), obtained
    under different protocols (the published rows are official test-set results over
    sequences 11--21 with the benchmark's own matching; this work is scored on
    sequence~08 by greedy point-level IoU$\,\ge\,$0.3 one-to-one matching, car class
    only), and over different scopes (all classes vs.\ car-only; online vs.\ offline).
    All published rows additionally require GPU training on dense per-point labels;
    this work's segmentation stage requires no training and no labels, and its
    classifier uses only automatically mined, weakly labelled clusters. The table
    compares \emph{requirements and paradigm}; on raw accuracy the supervised networks
    are stronger by design. Published figures are the test-set numbers reported by the
    cited papers; this work's row is the frozen sequence~08 headline (Chapter~\ref{ch:detection-eval},
    Finding~\#34). RangeNet++ appears in two rows: as a standalone semantic backbone
    (mIoU 52.2) and as the semantic half of the semantics-plus-detector panoptic
    baseline of~\cite{behley2021panoptic} (PQ 37.1).}
  \label{tab:positioning}
  \small
  \renewcommand{\arraystretch}{1.25}
  \setlength{\tabcolsep}{5pt}
  \begin{tabular}{@{}p{2.7cm}p{3.5cm}p{3.0cm}p{4.4cm}@{}}
    \toprule
    Method & Approach & Supervision & Reported on SemanticKITTI (own metric) \\
    \midrule
    RangeNet++~\cite{milioto2019rangenet++} & Range-image CNN, semantic & Dense per-point & mIoU 52.2; car IoU 91.4 \\
    RandLA-Net~\cite{Hu_2020_CVPR} & Point-based, semantic & Dense per-point & mIoU 53.9; car IoU 94.2 \\
    Cylinder3D~\cite{zhu2021cylindrical} & Voxel (cylindrical), semantic & Dense per-point & mIoU 67.8; car IoU 97.1 \\
    \addlinespace[2pt]
    RangeNet++ \& PointPillars~\cite{behley2021panoptic} & Semantic net + detector, panoptic & Dense per-point + boxes & PQ 37.1; PQ$^{\text{Th}}$ 20.2 \\
    Panoptic-PolarNet~\cite{zhou2021panopticpolarnet} & End-to-end panoptic & Dense per-point + instance & PQ 54.1; PQ$^{\text{Th}}$ 53.3 \\
    DS-Net~\cite{hong2021dsnet} & End-to-end panoptic & Dense per-point + instance & PQ 55.9; mIoU 61.6 \\
    EfficientLPS~\cite{sirohi2021efficientlps} & End-to-end panoptic & Dense per-point + instance & PQ 57.4; PQ$^\dagger$ 63.2 \\
    \midrule
    \textbf{This work} & Geometric clustering + binary classifier (+ completion) & None for segmentation; auto-mined weak labels for the classifier & car F$_1$ 0.808 (R 0.730, P 0.905); point-IoU$\,\ge\,$0.3, seq~08 \\
    \bottomrule
  \end{tabular}
\end{table}

Read as requirements, the table states the thesis's position plainly. The supervised
networks reach substantially higher accuracy on their own metrics, and this work does
not contest that; its car-only recall of \num{0.730} (against the
$\ge$-10-surviving-points denominator of Chapter~\ref{ch:detection-eval}) is well below the accuracy of those
networks --- a difference of paradigm, not a lower score on a shared axis, since none of
the published metrics is this work's point-level car F$_1$. What the pipeline offers instead is a different set of properties: it
segments without any labelled point clouds, it runs on a single consumer GPU without a
training stage for detection (its runtime is reported in Section~\ref{sec:results-runtime}), and every stage of
its decision is inspectable, so its
failure modes can be --- and are --- root-caused rather than left opaque (the recall
shortfall is traced to clustering fragmentation, not classifier error, in Chapter~\ref{ch:recall-bottleneck}).
The accuracy gap is therefore a deliberate consequence of choosing interpretability and
near-zero annotation over benchmark accuracy, not an unexplained deficiency.

A related question the table anticipates is why the comparison is drawn against
segmentation methods rather than 3D object detection, the task the ``detector'' label
most naturally invokes. Detection benchmarks score oriented 3D bounding boxes by average
precision, as in PointPillars~\cite{lang2019pointpillars} and
CenterPoint~\cite{yin2021centerpoint}. This pipeline, however, emits point-level clusters
--- it must, because those clusters are the input to the completion stage (Chapter~\ref{ch:completion-method}) ---
and never fits an oriented box; box AP would require a box-fitting stage the pipeline
deliberately does not have. Point-level IoU against the ground-truth car points is
therefore the metric that matches the actual output, and the choice of IoU threshold and
recall denominator that make that measurement precise --- together with its sensitivity
to those choices --- is examined directly in Section~\ref{sec:eval-protocol}.

This framing also sets the boundary of the thesis's contribution, which is not a claim
of pipeline novelty. Classical ground-removal-plus-clustering detection --- including recent
clustering-plus-classifier detectors for LiDAR vehicles~\cite{lim2025longrange} ---
a lightweight multi-object tracking stage used only to filter transient detections
across frames (in the tradition of 3D multi-object trackers such as
AB3DMOT~\cite{weng2020ab3dmot}, and not a contribution of this work), and PCN-style
completion are established building blocks, reimplemented and tuned here. The
contributions this work does claim are methodological and diagnostic: a leakage-free
metric for completion quality on real LiDAR where reference-based Chamfer distance is
invalid (Chapter~\ref{ch:completion-eval}), the root-cause analysis of the recall bottleneck together with the
chain of negative repair results (Chapter~\ref{ch:recall-bottleneck}), and the finding that synthetic
pretraining is redundant for the classifier at scale (Chapter~\ref{ch:detection-eval}). The positioning table
is the backdrop against which those contributions are argued: this is an interpretable,
annotation-light system whose value is in what it explains and measures, not in topping a
benchmark it was never designed to win.
